% Slug: ch23-inline-scope
% Source: chapters/ch23.tex:60 (tikzpicture env start line)
% Original label: \label{fig:ch23-inline-scope}
% Original caption (LaTeX source):
%   Inlining merges the callee body into the caller, so later passes see one region instead of two isolated ones.
% Original caption (plain text):
%   Inlining merges the callee body into the caller, so later passes see one region instead of two isolated ones.
% Required TikZ setup: \usepackage{tikz} + \usetikzlibrary{arrows.meta, positioning, shapes.geometric, fit} (loaded by support/compiler-pseudocode.sty).
% Shared styles used in this picture: `basic block', `cfg edge'.
%   Defined in support/compiler-pseudocode.sty: `basic block', `cfg edge', `back edge', `block label'.
% Usage: % Slug: ch23-inline-scope
% Source: chapters/ch23.tex:60 (tikzpicture env start line)
% Original label: \label{fig:ch23-inline-scope}
% Original caption (LaTeX source):
%   Inlining merges the callee body into the caller, so later passes see one region instead of two isolated ones.
% Original caption (plain text):
%   Inlining merges the callee body into the caller, so later passes see one region instead of two isolated ones.
% Required TikZ setup: \usepackage{tikz} + \usetikzlibrary{arrows.meta, positioning, shapes.geometric, fit} (loaded by support/compiler-pseudocode.sty).
% Shared styles used in this picture: `basic block', `cfg edge'.
%   Defined in support/compiler-pseudocode.sty: `basic block', `cfg edge', `back edge', `block label'.
% Usage: % Slug: ch23-inline-scope
% Source: chapters/ch23.tex:60 (tikzpicture env start line)
% Original label: \label{fig:ch23-inline-scope}
% Original caption (LaTeX source):
%   Inlining merges the callee body into the caller, so later passes see one region instead of two isolated ones.
% Original caption (plain text):
%   Inlining merges the callee body into the caller, so later passes see one region instead of two isolated ones.
% Required TikZ setup: \usepackage{tikz} + \usetikzlibrary{arrows.meta, positioning, shapes.geometric, fit} (loaded by support/compiler-pseudocode.sty).
% Shared styles used in this picture: `basic block', `cfg edge'.
%   Defined in support/compiler-pseudocode.sty: `basic block', `cfg edge', `back edge', `block label'.
% Usage: % Slug: ch23-inline-scope
% Source: chapters/ch23.tex:60 (tikzpicture env start line)
% Original label: \label{fig:ch23-inline-scope}
% Original caption (LaTeX source):
%   Inlining merges the callee body into the caller, so later passes see one region instead of two isolated ones.
% Original caption (plain text):
%   Inlining merges the callee body into the caller, so later passes see one region instead of two isolated ones.
% Required TikZ setup: \usepackage{tikz} + \usetikzlibrary{arrows.meta, positioning, shapes.geometric, fit} (loaded by support/compiler-pseudocode.sty).
% Shared styles used in this picture: `basic block', `cfg edge'.
%   Defined in support/compiler-pseudocode.sty: `basic block', `cfg edge', `back edge', `block label'.
% Usage: \input{graphs-code/ch23-inline-scope.tikz.tex} inside a figure environment
%   (wrap with \begin{figure} ... \caption{...} \label{fig:ch23-inline-scope} ... \end{figure}).
% Note: tikzpicture body below is verbatim from the source; do not reformat.
\begin{tikzpicture}[node distance=1.1cm and 1.6cm]
\node[basic block] (caller) {caller: t = call f(a)};
\node[basic block, right=of caller] (callee) {f: body; ret r};
\node[basic block, below=of caller] (merged) {caller: body[a/x]; use r};
\draw[cfg edge] (caller) -- (callee);
\draw[cfg edge] (caller) -- (merged);
\end{tikzpicture}
 inside a figure environment
%   (wrap with \begin{figure} ... \caption{...} \label{fig:ch23-inline-scope} ... \end{figure}).
% Note: tikzpicture body below is verbatim from the source; do not reformat.
\begin{tikzpicture}[node distance=1.1cm and 1.6cm]
\node[basic block] (caller) {caller: t = call f(a)};
\node[basic block, right=of caller] (callee) {f: body; ret r};
\node[basic block, below=of caller] (merged) {caller: body[a/x]; use r};
\draw[cfg edge] (caller) -- (callee);
\draw[cfg edge] (caller) -- (merged);
\end{tikzpicture}
 inside a figure environment
%   (wrap with \begin{figure} ... \caption{...} \label{fig:ch23-inline-scope} ... \end{figure}).
% Note: tikzpicture body below is verbatim from the source; do not reformat.
\begin{tikzpicture}[node distance=1.1cm and 1.6cm]
\node[basic block] (caller) {caller: t = call f(a)};
\node[basic block, right=of caller] (callee) {f: body; ret r};
\node[basic block, below=of caller] (merged) {caller: body[a/x]; use r};
\draw[cfg edge] (caller) -- (callee);
\draw[cfg edge] (caller) -- (merged);
\end{tikzpicture}
 inside a figure environment
%   (wrap with \begin{figure} ... \caption{...} \label{fig:ch23-inline-scope} ... \end{figure}).
% Note: tikzpicture body below is verbatim from the source; do not reformat.
\begin{tikzpicture}[node distance=1.1cm and 1.6cm]
\node[basic block] (caller) {caller: t = call f(a)};
\node[basic block, right=of caller] (callee) {f: body; ret r};
\node[basic block, below=of caller] (merged) {caller: body[a/x]; use r};
\draw[cfg edge] (caller) -- (callee);
\draw[cfg edge] (caller) -- (merged);
\end{tikzpicture}
